\section{Answers to review questions}
\label{sec:questions}

Twelve questions were raised on an earlier draft. Six can be answered from the runs we
already have; the rest are proposed work, with what each depends on.

\subsection{Answered from the data we have}

\paragraph{Why are there more components than services?}
One service does run in one container, but a deployment also contains databases, a
message broker, a cache, a service registry and a proxy. The small shop is 14
containers: 8 business services, 4 databases, a broker and an ingress. The large
booking system is about 67 containers: roughly 39 services plus a shared database, a
second database, a cache, a registry and a dashboard. The extra components matter
because they are exactly the ones that emit no request traces --- they appear only as
something a caller waited on.

\paragraph{Are faults in services found more easily than faults in other components?}
Yes, and by a wide margin.

\begin{table}[h]
\centering\small
\begin{tabular}{lccc}
\toprule
Where the fault was injected & Cases & Right component & Fully correct \\
\midrule
The whole machine            & 8  & 7/8  & 6/8 \\
A business service           & 10 & 9/10 & 4/10 \\
A non-service component      & 3  & 1/3  & 1/3 \\
A rogue container            & 2  & 2/2  & 0/2 \\
\bottomrule
\end{tabular}
\end{table}

\noindent
Faults in services are localized 9 times in 10; faults in databases and queues only 1
in 3. This is the blind spot we predicted before running anything: those components
produce no traces, so ordinary telemetry points \emph{near} them, never \emph{at}
them. The one success in that row is the slow database, and it is carried by the
kernel wait analysis --- it stays correct even with traces removed entirely.

\paragraph{What does the per-fault table look like?}
It is already in the appendix: 23 faults $\times$ 7 setups, generated from the
recorded results.

\paragraph{Are the investigation guides specific to one root cause?}
Indirectly, yes. When we hide the matching guide, the selector picks a wrong one and
that costs about 18 points versus using no guide at all. A wrong guide actively
misleads, which only happens if guides are specific.

\paragraph{How do the other methods behave when data is reduced?}
We measured this: the statistical method falls from 38\% to 34\% as traces are
thinned; the academic multi-modal method stays flat at 48\%. Flat sounds robust but
is not --- each leans on a single data source and is already low, and neither can use
the kernel layer at all. Feeding kernel data into the academic method changed nothing.

\paragraph{Is the agent only better because it sees more data?}
No, and we can show it. Given \emph{115 times less data and no kernel layer at all},
the agent still names the right component 83\% of the time. The other methods, given
\emph{everything}, reach 46--48\%. The advantage survives when the data advantage is
reversed, so what is left is the iterative questioning, not the volume.

\subsection{Proposed next, on data we already hold}

\paragraph{A language-model-only baseline.} Same model, no tools, one shot: it sees
only the evidence briefing and must answer in the same format. This closes the
question of whether the result comes from the agent loop or from the model alone. We
will follow a published prompt design for one-shot incident diagnosis so the
comparison is defensible. Under \$1.

\paragraph{A guide cross-matrix.} Force every guide onto every fault instead of
letting the selector choose. If guides are specific we should see a strong diagonal
and a depressed off-diagonal. That single figure is the specificity claim. About \$5.

\paragraph{Ranked answers and ranking metrics.} The agent currently gives one verdict.
We will let it return up to five candidates, each carrying its own evidence, and
report hit rate at 1/3/5 along with mean reciprocal rank, precision and average
precision --- directly comparable to the other methods, which already produce ranked
lists. One caveat we will state plainly: a language model's list is ordered by what
its reasoning found plausible next, not by calibrated probability. We reduce that by
dropping any candidate without its own evidence, and can compare against a variant
that ranks by how often repeated runs agree. Roughly \$1--2.

\subsection{Proposed next, but blocked on new recordings}

\paragraph{Faults injected into source code.} The agent has a code-search tool that
today's faults never exercise, because all twelve are operational. The mutations have
to be ones that actually show up in telemetry --- removing a downstream call, removing
a logged statement, changing a database query --- not something invisible like
flipping a condition. Our two instrumented services can be rebuilt from a branch, so
the mechanism is ready.

\paragraph{More than one fault at a time.} Every recording currently has exactly one
injected fault, which is what makes the labels unambiguous. Adding pairs --- some
overlapping in symptoms, some independent --- would need the set-aware metrics above,
because three guesses covering two true faults is not the same as three guesses for
one.

\noindent
Both need the collection machine, which is currently switched off. Everything up to
that point can be prepared now.

\subsection{Noted, not pursued}

Mapping data sources to fault types is a related but different direction from ours
(we map fault types to investigation guides), and is already being published
elsewhere; we will cite it rather than enter it. A pipeline diagram of the whole
system, and a table showing that each guide encodes a signature written down before
any experiment, are both write-up work for the paper version.
